\noindent
\textbf{\large LISTA DE ACRÓNIMOS Y SÍMBOLOS}

\vspace{1cm}

\noindent\textbf{Infraestructura y arquitectura de software}

\vspace{0.3cm}

\begin{description}[leftmargin=3cm, style=sameline, font=\normalfont]
  \item[API]     \textit{Application Programming Interface}; interfaz de programación que
                 expone funcionalidades de un servicio para ser consumidas por otros módulos
                 o clientes externos.
  \item[REST]    \textit{Representational State Transfer}; estilo arquitectónico para el
                 diseño de APIs HTTP basado en recursos y verbos estándar.
  \item[HTTP /]  \textit{HyperText Transfer Protocol (Secure)}; protocolos de transporte
  \item[HTTPS]   sobre los que se exponen las APIs del sistema.
  \item[JSON]    \textit{JavaScript Object Notation}; formato de intercambio de datos
                 estructurados utilizado en los mensajes entre módulos.
  \item[URL]     \textit{Uniform Resource Locator}; identificador de recursos web utilizado
                 para apuntar a endpoints y nodos del sistema.
  \item[RPC]     \textit{Remote Procedure Call}; modelo de invocación remota empleado en la
                 comunicación con nodos de la red BlockDAG.
  \item[SQL]     \textit{Structured Query Language}; lenguaje de consulta utilizado en la
                 base de datos relacional de metadatos operativos.
  \item[CPU]     \textit{Central Processing Unit}; procesador principal del host, recurso
                 monitoreado en los experimentos de rendimiento.
  \item[RAM]     \textit{Random Access Memory}; memoria principal del host, recurso
                 monitoreado en los experimentos de escalabilidad.
  \item[IoT]     \textit{Internet of Things}; paradigma de dispositivos interconectados;
                 fuente de datos real a integrar en despliegues futuros.
  \item[FastAPI] \textit{FastAPI}; framework web asíncrono de Python utilizado para
                 implementar las APIs REST de los módulos del sistema.
\end{description}

\vspace{0.8cm}
\noindent\textbf{Simulación y tráfico urbano}

\vspace{0.3cm}

\begin{description}[leftmargin=3cm, style=sameline, font=\normalfont]
  \item[SUMO]    \textit{Simulation of Urban MObility}; plataforma de simulación de tráfico
                 urbano de código abierto utilizada en el módulo \textit{traffic-sim}.
  \item[TraCI]   \textit{Traffic Control Interface}; protocolo de comunicación de SUMO que
                 permite controlar y consultar la simulación en tiempo real.
  \item[ITS]     \textit{Intelligent Transportation Systems}; sistemas inteligentes de
                 transporte; marco general en el que se inscribe este trabajo.
  \item[MILP]    \textit{Mixed-Integer Linear Programming}; programación entera mixta;
                 formulación clásica del problema de optimización de planes semafóricos.
  \item[TPS]     Transacciones por segundo; métrica de rendimiento empleada para comparar
                 el throughput de BlockDAG frente a blockchain lineal.
\end{description}

\vspace{0.8cm}
\noindent\textbf{Inteligencia computacional}

\vspace{0.3cm}

\begin{description}[leftmargin=3cm, style=sameline, font=\normalfont]
  \item[PSO]     \textit{Particle Swarm Optimization}; algoritmo de optimización por
                 enjambre de partículas utilizado para ajustar los tiempos semafóricos.
  \item[FIS]     \textit{Fuzzy Inference System}; sistema de inferencia difusa; implementado
                 con modelo Mamdani para clasificar el nivel de congestión.
\end{description}

\vspace{0.8cm}
\noindent\textbf{Almacenamiento descentralizado}

\vspace{0.3cm}

\begin{description}[leftmargin=3cm, style=sameline, font=\normalfont]
  \item[IPFS]      \textit{InterPlanetary File System}; sistema de almacenamiento distribuido
                   por contenido utilizado para persistir los datos de tráfico.
  \item[CID]       \textit{Content Identifier}; identificador criptográfico de contenido en
                   IPFS; permite verificar la integridad de los datos almacenados.
  \item[DAG]       \textit{Directed Acyclic Graph}; grafo dirigido acíclico; estructura de
                   datos base de IPFS y de las redes BlockDAG.
  \item[BlockDAG]  Red de bloques organizada como grafo dirigido acíclico en lugar de cadena
                   lineal; permite mayor throughput y confirmación rápida de transacciones.
  \item[PHANTOM /] Protocolos de consenso y ordenamiento para redes BlockDAG que definen
  \item[GHOSTDAG]  la cadena principal y resuelven conflictos entre bloques paralelos.
  \item[PoW]       \textit{Proof of Work}; mecanismo de consenso basado en cómputo
                   intensivo; referencia comparativa frente a BlockDAG.
  \item[PoS]       \textit{Proof of Stake}; mecanismo de consenso basado en participación
                   económica; alternativa energéticamente eficiente a PoW.
  \item[P2P]       \textit{Peer-to-Peer}; modelo de red entre pares utilizado por IPFS y las
                   redes BlockDAG para distribuir contenido y bloques.
  \item[DHT]       \textit{Distributed Hash Table}; tabla de hash distribuida que usa IPFS
                   para localizar nodos que almacenan un CID dado.
  \item[Solidity]  Lenguaje de programación de contratos inteligentes utilizado para
                   implementar el registro inmutable en la red BlockDAG.
  \item[SHA]       \textit{Secure Hash Algorithm}; familia de funciones de resumen
                   criptográfico empleadas en la construcción de CIDs e identificadores de
                   bloque.
\end{description}

\vspace{0.8cm}
\noindent\textbf{Símbolos matemáticos}

\vspace{0.3cm}

\begin{description}[leftmargin=3cm, style=sameline, font=\normalfont]
  \item[$x_i$]       Posición de la partícula $i$ en el espacio de búsqueda del PSO;
                     representa una configuración candidata de tiempos semafóricos.
  \item[$v_i$]       Velocidad de la partícula $i$ en el PSO; determina la magnitud y
                     dirección del desplazamiento en cada iteración.
  \item[$p_i$]       Mejor posición histórica (\textit{personal best}) de la partícula $i$.
  \item[$g$]         Mejor posición global (\textit{global best}) encontrada por el enjambre.
  \item[$w$]         Factor de inercia del PSO; controla el peso de la velocidad previa en
                     la actualización de la trayectoria.
  \item[$c_1, c_2$]  Coeficientes de aceleración cognitivo y social del PSO.
  \item[$r_1, r_2$]  Números aleatorios uniformes en $[0,1]$ usados en la ecuación de
                     actualización de velocidad del PSO.
  \item[$C$]         Longitud de ciclo semafórico; duración total de una repetición completa
                     de las fases.
  \item[$g_k$]       Tiempo de verde asignado a la fase $k$ dentro del ciclo $C$.
  \item[$\mu$]       Grado de pertenencia en lógica difusa; valor en $[0,1]$ que indica
                     cuánto pertenece un valor a un conjunto difuso.
  \item[$\phi_j$]    Función de pertenencia $j$; define la forma del conjunto difuso
                     (triangular, trapezoidal, gaussiana).
\end{description}
